\chapter{Layer 9: WORM Provenance Receipts}
\label{ch:worm}

\section{Purpose and Contract}
\label{sec:worm-purpose}

Layer~9 is the historical memory of the pipeline. Its inbound contract is
``an artifact produced by any upstream layer''; its outbound contract is ``a
signed, hash-linked receipt appended to an immutable ledger.'' Without this
layer, every earlier guarantee would be untraceable: we could know a result
is correct but not \emph{by whom, from what, and when}.

\section{The Receipt}
\label{sec:worm-receipt}

The reference receipt generator (\texttt{mathlib5-ffi-bridge/scripts/receipt.sh})
hashes every source file and the built binary, then writes a manifest and a
seal:
\begin{lstlisting}[language=BashX,caption={Receipt generation (receipt.sh).},label=lst:worm-receipt]
find "$BRIDGE_DIR/C" "$BRIDGE_DIR/Lean" -type f \( \
     -name '*.hs' -o -name '*.lean' -o -name '*.h' -o -name '*.c' \) | \
  sort | xargs sha256sum > "$RECEIPT_DIR/source.sha256"

cat > "$RECEIPT_DIR/manifest.json" << EOF
{
  "build_id": "receipt-$(date +%s)",
  "timestamp": "$(date -u +%Y-%m-%dT%H:%M:%SZ)",
  "git_sha": "$(git -C "$BRIDGE_DIR" rev-parse HEAD 2>/dev/null || echo unknown)",
  "lean_toolchain": "$(cat "$BRIDGE_DIR/lean-toolchain" 2>/dev/null || echo unknown)",
  "source_files": $(wc -l < "$RECEIPT_DIR/source.sha256"),
  "status": "BUILT"
}
EOF
sha256sum "$RECEIPT_DIR/manifest.json" | cut -d' ' -f1 > "$RECEIPT_DIR/seal.sha256"
\end{lstlisting}
Each receipt therefore binds: the exact source hashes, the build identity, the
git commit, the toolchain, and a final seal. The seal is itself a SHA-256,
making the receipt self-verifying.

\section{The Metadata Schema}
\label{sec:worm-meta}

Every component carries a \texttt{metadata.json} (see \cref{app:metadata} for
the full schema). Its security-audit block is uniform across the repository:
\begin{lstlisting}[language=,caption={Uniform security-audit block (metadata.json).},label=lst:worm-meta]
"security_audit": {
    "metadata_integrity": "Yes",
    "tamper_evidence":    "SHA-256 chain (append-only)",
    "encryption":         "AES-256-GCM",
    "plasma_gate":        "Ed25519_Enforced"
}
\end{lstlisting}
The \texttt{status} field is almost always \texttt{awaiting\_worm\_seal},
recording that the artifact is staged but not yet finally sealed into the
Bifrost chain. The \texttt{provenance} block names the
\texttt{chain\_link} (\texttt{Bifrost\_WORM\_Chain\_20260710\_01}), the
\texttt{source\_path}, and a \texttt{corpora\_path} --- a three-way pointer
that locates the artifact in source, in the ledger, and in the training
corpus.

\section{Hash-Chain Tamper Evidence}
\label{sec:worm-chain}

Receipts are appended to the Bifrost WORM chain. Let $t_i$ be receipt $i$.
Then
\begin{equation}
t_i = H\bigl(t_{i-1} \;\Vert\; \text{payload}_i\bigr),
\end{equation}
so the ledger is a hash chain. An attacker who rewrites an early receipt must
recompute every subsequent link; without the private key to re-sign (Plasma
Gate, \cref{sec:coh-plasma}), the forgery is detectable. This is the
``immutable system'' of the title made concrete.

\section{Why Provenance Is a Layer, Not a Log}
\label{sec:worm-why}

Provenance is placed \emph{inside} the pipeline rather than bolted on as an
external log because every other layer depends on it: policy (L7) consults
receipts; the swarm (L8) submits sealed solutions; completeness (L10) records
which proofs were closed by whom. Making provenance a first-class boundary
means no artifact can bypass it.
